\documentclass[12pt]{article}
\usepackage[utf8]{inputenc}
\usepackage[T1]{fontenc}
\usepackage{hyperref}
\title{Tesaka Malagasy Language Resources for K-12 Curriculum Development}
\author{Compiled by Claude Sonnet 5 (claude-sonnet-5) for Andr\'as Kornai}
\date{}
\begin{document}
\maketitle

\section{Introduction}
Tesaka (also Antesaka or Antaisaka) is a Malagasy variety spoken by the Antesaka people in the Atsimo Atsinanana region of southeastern Madagascar. Searches found no dedicated grammar, dictionary, textbook or corpus for Tesaka itself. Sources also disagree on its code and speaker numbers (one gives tkg with about 1.97 million speakers, another about 1.13 million), so these should be checked against SIL or Glottolog.

\section{Primary Educational Resources}
No Tesaka-specific primers, textbooks or school programmes were found in the searches. Official Malagasy (Merina-based) materials are the likely basis for any curriculum, with Tesaka adaptations to be developed.

\section{Linguistic and Reference Materials}
\begin{itemize}
\item Antesaka dialect (Wikipedia overview, with code and region): \url{https://en.wikipedia.org/wiki/Antesaka_dialect}
\item Academic paper on Malagasy dialect classification and dictionary coverage, which notes that documentation is incomplete for many dialects: \url{https://www.academia.edu/8193206}
\item Antaisaka (langue), a French-language overview with sample lexical differences from Merina: \url{https://franco.wiki/fr/Antaisaka_%28langue%29.html}
\end{itemize}

\section{Religious and Community Resources}
None found in the searches.

\section{Digital and Audio Resources}
The Wikipedia Antesaka article embeds recorded speech from Midongy Atsimo and Efatsy speakers: \url{https://en.wikipedia.org/wiki/Antesaka_dialect}

\section{Curriculum Development Context}
Very little Tesaka-specific material was found. A curriculum would need new fieldwork and community-created materials, built on the standard Malagasy resources. Further searches under the names Antaisaka and Antesaka are recommended.
\end{document}